\subsection{Кошмар в замке}

\subsubsection{Описание решения}

По условию задачи, на стоимость строки влияет каждая пара уникальных одинаковых символов. Если они встречаются чаще, чем два раза --- третий и далее символы не вносят свой вклад, их можно размещать как угодно.

Значит, при вводе исходной строки стоит проанализировать её и подсчитать статистику по каждому символу: сколько раз он встречается в строке. Если он появился хотя бы два раза, его можно использовать для повышения стоимости строки.

Далее в задаче говорится, что стоимость считается как сумма произведений веса пары символов на максимальное расстояние между ними. Поступим жадно: пару символов, которая весит больше, разместим на большее расстояние. Чтобы этого достичь, будем размещать их по принципу палиндрома: \( i \)-ую по весу пару символов разместим на \( i \)-ое и \( N-i \)-ое места. Оставшиеся символы добавим в центр в произвольном порядке --- на стоимость строки это не повлияет.

\subsubsection{Асимптотическая сложность}

\textit{Временная сложность решения} --- линейная \( \mathcal{O}(N) \). В алгоритме самая затратная по времени операция --- посимвольная обработка строки: на каждый символ отводится константное число операций.

Полезная работа по максимизации стоимости строки выполняется за константное время, поскольку зависит от алфавита, который используется во входной строке --- последний фиксирован по условию задачи.

\textit{Пространственная сложность решения} --- константная \( \mathcal{O}(1) \). Входная строка считывается посимвольно и не сохраняется в памяти целиком. Алфавит имеет фиксированный размер, поэтому на его хранение тоже требуется константное число ячеек.

\subsubsection{Листинг кода}

\begin{lstlisting}
#include <algorithm>
#include <iostream>
#include <unordered_map>
#include <vector>

using letter = struct {
  char c;
  long v;
};

int main() {
  std::unordered_map<char, long> occur;
  char c;
  while (std::cin >> std::noskipws >> c && c != '\n') {
    occur[c] = occur.count(c) ? occur[c] + 1 : 1;
  }
  std::cin >> std::skipws;
  std::vector<letter> costs;
  struct {
    bool operator()(letter a, letter b) const {
      return a.v > b.v;
    }
  } customLess;
  for (int i = 0; i < 26; i++) {
    letter l;
    l.c = 'a' + i;
    std::cin >> l.v;
    costs.push_back(l);
  }
  std::sort(costs.begin(), costs.end(), customLess);
  std::string pair, rest;
  for (letter l : costs) {
    if (occur.count(l.c) == 0) {
      continue;
    }
    if (occur[l.c] > 1) {
      pair += l.c;
      for (int i = 0; i < occur[l.c] - 2; i++) {
        rest += l.c;
      }
    } else {
      rest += l.c;
    }
  }
  std::cout << pair << rest << std::string(pair.rbegin(), pair.rend());
  return 0;
}
\end{lstlisting}